\subsection{Resumen de Calendarización y Presupuesto - El tiempo y los recursos detrás de la entrega}\label{sec:spmp-6-5}
El periodo 2027-1 comprende Construcción y Semi-transición según el ciclo de vida de la sección~\ref{sec:spmp-8-1}; Inicio en 2026-3 aporta la preparación del proyecto. Los productos se describen en la sección~\ref{sec:spmp-6-4}~\cite{adaptacion2026}. La tabla~\ref{tab:presupuesto} resume el presupuesto; el detalle mensual y el flujo de caja se presentan en la tabla~\ref{tab:flujo-caja}. El trabajo del equipo se valora como aporte en especie, no como pago.

\begin{longtable}{|>{\raggedright\arraybackslash}p{0.13\dimexpr\linewidth-10\tabcolsep-6\arrayrulewidth\relax}|>{\raggedright\arraybackslash}p{0.23\dimexpr\linewidth-10\tabcolsep-6\arrayrulewidth\relax}|>{\raggedright\arraybackslash}p{0.06\dimexpr\linewidth-10\tabcolsep-6\arrayrulewidth\relax}|>{\raggedright\arraybackslash}p{0.44\dimexpr\linewidth-10\tabcolsep-6\arrayrulewidth\relax}|>{\raggedright\arraybackslash}p{0.14\dimexpr\linewidth-10\tabcolsep-6\arrayrulewidth\relax}|}
\caption{Marco temporal de referencia. Fuente: metodología del equipo~\cite{adaptacion2026}.}\label{tab:hitos}\\\hline
\textbf{Periodo} & \textbf{Fase y orientación} & \textbf{\textcolor{black}{ID}} & \textbf{\textcolor{black}{Hito}} & \textbf{\textcolor{black}{Fecha}}\\\hline
\endfirsthead\hline
\textbf{Periodo} & \textbf{Fase y orientación} & \textbf{\textcolor{black}{ID}} & \textbf{\textcolor{black}{Hito}} & \textbf{\textcolor{black}{Fecha}}\\\hline
\endhead\multirow{2}{=}{2026-3} & \multirow[t]{2}{=}{Inicio: preparación del proyecto.} & \textcolor{black}{H1} & \textcolor{black}{Haber aprendido a usar FDS y sus capacidades para incendios forestales (WFDS).} & \textcolor{black}{28/10/2026} \\*\cline{3-5}\noalign{\nobreak} & & \rule[-3\baselineskip]{0pt}{3.8\baselineskip}\textcolor{black}{H2} & \textcolor{black}{Haber organizado las fuentes de datos necesarias para las simulaciones.} & \textcolor{black}{11/11/2026} \\\hline\multirow{2}{=}{2027-1} & \multirow[t]{2}{=}{Construcción: desarrollo e integración de la aplicación y \textcolor{black}{del dataset}.} & \textcolor{black}{H3} & \textcolor{black}{Haber generado \textcolor{black}{el dataset completo de simulaciones}.} & \textcolor{black}{02/04/2027} \\*\cline{3-5}\noalign{\nobreak} & & \rule[-3\baselineskip]{0pt}{3.8\baselineskip}\textcolor{black}{H4} & \textcolor{black}{Haber desarrollado la aplicación de escritorio.} & \textcolor{black}{09/04/2027} \\\hline\noalign{\break}\multirow{2}{=}{2027-1} & \multirow[t]{2}{=}{Semi-transición: evaluación, correcciones y entrega académica.} & \textcolor{black}{H5} & \textcolor{black}{Haber evaluado el producto con usuarios.} & \textcolor{black}{23/04/2027} \\*\cline{3-5}\noalign{\nobreak} & & \rule[-3\baselineskip]{0pt}{3.8\baselineskip}\textcolor{black}{H6} & \textcolor{black}{Haber consolidado y revisado con el director el paquete final del proyecto.} & \textcolor{black}{07/05/2027} \\\hline
\end{longtable}

\begin{longtable}{|>{\raggedright\arraybackslash}p{0.21\dimexpr\linewidth-6\tabcolsep-4\arrayrulewidth\relax}|>{\raggedright\arraybackslash}p{0.55\dimexpr\linewidth-6\tabcolsep-4\arrayrulewidth\relax}|>{\raggedright\arraybackslash}p{0.24\dimexpr\linewidth-6\tabcolsep-4\arrayrulewidth\relax}|}
\caption{Ítems consolidados del presupuesto. Fuente: estimación del equipo.}\label{tab:presupuesto}\\
\hline
\textbf{Ítem de presupuesto} & \textbf{Descripción} & \textbf{Costo asociado (COP)}\tabularnewline\hline
\endfirsthead
\hline
\textbf{Ítem de presupuesto} & \textbf{Descripción} & \textbf{Costo asociado (COP)}\tabularnewline\hline
\endhead
\raggedright Trabajo del equipo & \raggedright 1.046 horas-persona a \$27.097 por hora, a precios de 2026 y sin cargas prestacionales: 384 horas de capacidad formal en Inicio (cuatro integrantes, dos créditos y 48 horas por crédito), valoradas en \$10.405.248, y 662 horas estimadas para 2027-1, valoradas en \$17.938.214. & \raggedright \$28.343.462\tabularnewline\hline
\raggedright Infraestructura HPC-ZINE & \raggedright \textcolor{black}{Reserva provisional de electricidad para A5 y A6, calculada con un tope de 10.000 núcleo-horas. La tarifa del servicio y la forma de cobertura se confirmarán con HPC-ZINE (sección~\ref{sec:spmp-9-3-4}).} & \raggedright \textcolor{black}{\$280.000}\tabularnewline\hline
\raggedright Transporte & \raggedright \$100.000 para la sesión de entrevistas con usuarios en 2026-3 y \$300.000 para tres presentaciones a usuarios en \textcolor{black}{abril de 2027} (A9). & \raggedright \$400.000\tabularnewline\hline
\raggedright Software y licencias & \raggedright Las herramientas de la sección~8.2 son gratuitas. & \raggedright \$0\tabularnewline\hline
\raggedright \textbf{\textcolor{black}{Total económico estimado}} & \raggedright \textcolor{black}{\$400.000 en salidas de dinero, \$28.343.462 en trabajo del equipo y \$280.000 de reserva provisional para HPC-ZINE.} & \raggedright \textbf{\textcolor{black}{\$29.023.462}}\tabularnewline\hline
\end{longtable}

La tarifa de valoración se obtiene de 3,25~SMMLV de 2026 (\$1.750.905), punto medio del rango del Observatorio Laboral para un ingeniero de sistemas recién graduado, dividido entre 210 horas mensuales (sección~\ref{sec:spmp-9-3-4}). También se aportan en especie, sin valoración monetaria, el tiempo del director y del asesor, los portátiles de los integrantes y el servidor documental.
